\subsection{指标理论与 Lagrangian Grassmann 流形}\label{sec:11e}
设 $V=(V,\allowbreak \omega_V)$ 是带辛结构的 $2n$ 维向量空间，$\operatorname{Gr}(V)\cong U_n/O_n$ 是其无定向线性 Lagrangian 子空间的 Grassmann 流形。映射 $B\to\operatorname{Gr}(V)$ 等价于平凡辛向量丛 $B\times V\to B$ 的一个 Lagrangian 子丛，因此可以将 $\operatorname{Gr}(V)$ 的上同调类视为相对示性类。我们关心的是 Maslov 类 $\mu\in H^1(\operatorname{Gr}(V))$，以及由所有无定向子空间的通常 Grassmann 流形诱导的第二 Stiefel–Whitney 类 $w_2\in H^2(\operatorname{Gr}(V);\Z/2)$；第一 Stiefel–Whitney 类是 $\mu$ 的模 $2$ 约化，因此不必另行考虑。与 $\mu$ 配对给出同构 $\pi_1(\operatorname{Gr}(V))\to\Z$。当 $n\geq3$ 时，$w_2$ 同样给出同构 $\pi_2(\operatorname{Gr}(V))\to\Z/2$。低维例外是：$n=1$ 时 $U_1/O_1=S^1$，所以 $\pi_2(\operatorname{Gr}(V))=0$；$n=2$ 时 $U_2/O_2=S^1\times_{\Z/2}S^2$，所以 $\pi_2(\operatorname{Gr}(V))\cong\Z$。最后这个群可由二重覆盖 $U_2/SO_2$ 上的 Euler 类检测，而该类模 $2$ 约化为第二 Stiefel–Whitney 类，所以 $w_2:\pi_2(U_2/O_2)\to\Z/2$ 仍为满射。\znote{158}{Maslov 类是 Lagrangian Grassmann 流形的基本上同调类。nLab 指出 $H^1(\mathrm{LGrass},\Z)\cong\Z$，其生成元称为万有 Maslov 类 $u\in H^1(\mathrm{LGrass},\Z)$；对 Lagrangian 子流形 $Y$，把切丛的分类映射 $i:Y\to\mathrm{LGrass}$ 拉回，就得到 $i^*u\in H^1(Y,\Z)$，即 $Y$ 的 Maslov 类。这一上同调解释出自 Arnol'd 1967 年的工作，它把 Maslov 原来的渐近分析不变量提升为万有示性类。正文中 $\mu\in H^1(\operatorname{Gr}(V))$ 正是该类的具体化，因此与 $\mu$ 配对给出同构 $\pi_1(\operatorname{Gr}(V))\cong\Z$。对一条环路而言，配对得到的整数就是这条环路的 Maslov 指标：它把环路对应的一族 Lagrangian 子空间与一个固定参考子空间作交，再按穿越方向带符号求和；所得整数模 $2$ 约化即第一 Stiefel–Whitney 类，故正文说后者不必另行考虑。[来源：nLab "Maslov index"，https:/\allowbreak /\allowbreak ncatlab.org/\allowbreak nlab/\allowbreak show/\allowbreak Maslov+\allowbreak index（2026-11-15 访问）]}

还需要考察带基点环路空间 $\Omega\operatorname{Gr}(V)$ 与自由环路空间 $\mathcal L\operatorname{Gr}(V)$。它们各有以 $\Z$ 编号的连通分支 $\Omega_k\operatorname{Gr}(V)$ 或 $\mathcal L_k\operatorname{Gr}(V)$，由 Maslov 指标 $\mu(\rho)=\langle\mu,\allowbreak [\rho]\rangle=k$ 区分。当 $n\geq3$ 时，容易看出 $\pi_1(\Omega_k\operatorname{Gr}(V))\cong\Z/2$，并且
\[
\pi_1(\mathcal L_k\operatorname{Gr}(V))\cong\pi_2(\operatorname{Gr}(V))\times\pi_1(\operatorname{Gr}(V))\cong\Z/2\times\Z.
\]
特别地，
\begin{equation}\label{eq:11.15}
H^1(\mathcal L_k\operatorname{Gr}(V);\Z/2)\cong\Z/2\oplus\Z/2,
\end{equation}
其显式生成元为 $T(w_2)$ 与 $U(\mu)$ 的模 $2$ 约化。这里 $(T,\allowbreak U)$ 表示拉回映射
\[
H^*(\operatorname{Gr}(V))\longrightarrow H^*(S^1\times\mathcal L\operatorname{Gr}(V))
\cong H^{*-1}(\mathcal L\operatorname{Gr}(V))\oplus H^*(\mathcal L\operatorname{Gr}(V))
\]
的两个分量。

为与指标理论联系，作如下构造。固定 $V$ 上与 $\omega_V$ 相容的复结构 $I_V$、一个只有一个边界分支的紧 Riemann 曲面 $S$，以及保持定向的识别 $\partial S\cong S^1$。对每个 $\rho\in\mathcal L\operatorname{Gr}(V)$，令 $D_{S,\allowbreak \rho}$ 为作用于平凡复向量丛 $E=S\times V$ 的截面的 $\bar\partial$ 算子，边界条件由 Lagrangian 子丛 $F_\rho\subset E|\partial S$ 给出，$(F_\rho)_{\exp(2\pi is)}=\rho(s)$。考虑随 $\rho$ 变化的算子族 $D_{S,\allowbreak \rho}$，
\SourcePage{166}
便得到行列式线丛 $\underline{\det}_S\to\mathcal L\operatorname{Gr}(V)$。
\begin{lemma}\label{st:11.7}
$\operatorname{index}(D_{S,\allowbreak \rho})=n\chi(S)+\mu(\rho)$。此外，行列式线丛的同构类由下式给出：
\begin{equation}\label{eq:11.16}
w_1(\underline{\det}_S)=T(w_2)+(T(\mu)-1)U(\mu).
\end{equation}
\end{lemma}
\begin{proof}
第一条是熟知的带边界曲面的 Riemann–Roch 定理。第二条同样可视为“Real”算子族的一般指标定理\bcite{9}的特例。另一种论证不那么直接，却更为显式，它依赖以下几个要点。

\noindent(i) \emph{沿边界捏合分离。}假设 $S$ 的亏格大于零。可将它视为圆盘 $D$ 与同亏格闭 Riemann 曲面的连通和。闭曲面上的 $\bar\partial$ 算子显然是 $\CC$ 线性的，故其行列式线具有典范平凡化。于是由\eref{11.12}得到同构 $\underline{\det}_S\cong\underline{\det}_D$。

\noindent(ii) \emph{直和公式。}令 $\Delta:\operatorname{Gr}(V)\times\operatorname{Gr}(W)\to\operatorname{Gr}(V\oplus W)$ 为直和运算 $(\Lambda_1,\allowbreak \Lambda_2)\mapsto\Lambda_1\oplus\Lambda_2$，令
\[
\mathcal L_{k,l}\Delta:\mathcal L_k\operatorname{Gr}(V)\times\mathcal L_l\operatorname{Gr}(W)\longrightarrow\mathcal L_{k+l}\operatorname{Gr}(V\oplus W)
\]
为自由环路空间上的诱导映射。由\eref{11.1}，
\[
(\mathcal L_{k,l}\Delta)^*\underline{\det}_{S,V\oplus W}\cong\underline{\det}_{S,V}\boxtimes\underline{\det}_{S,W}.
\]
这里在记号中加入了所涉及的向量空间，以明确彼此关系。另一方面，运用 Whitney 和公式，
\[
\begin{aligned}
(\mathcal L_{k,l}\Delta)^*T(w_2)
&=T(w_2\times1+1\times w_2+\mu\times\mu)\\
&=T(w_2)\times1+1\times T(w_2)+U(\mu)\times T(\mu)+T(\mu)\times U(\mu)\\
&=(T(w_2)+lU(\mu))\times1+1\times(T(w_2)+kU(\mu)),\\
(\mathcal L_{k,l}\Delta)^*U(\mu)
&=U(\Delta^*\mu)=U(\mu)\times1+1\times U(\mu).
\end{aligned}
\]
因此
\[
\begin{aligned}
&(\mathcal L_{k,l}\Delta)^*(T(w_2)+(k+l-1)U(\mu))\\
&\quad=(T(w_2)+(k-1)U(\mu))\times1+1\times(T(w_2)+(l-1)U(\mu)).
\end{aligned}
\]
结合\eref{11.15}中生成元的表达式，上述计算还表明，只要 $V,\allowbreak W$ 中至少一个空间的维数不小于 $6$，$\mathcal L_{k,\allowbreak l}\Delta$ 在 $H^1(-;\Z/2)$ 上诱导的映射就是单射；事实上维数不小于 $4$ 已经足够，不过这里不强调这一点。

\noindent(iii) \emph{$n=1$ 时的计算。}假设 $n=1$，此时 Maslov 数 $\mu(\rho)$ 就是\secref{11d}考虑的绕数。例如，若 $S=D$ 且 $\mu(\rho)=-1$，则由\stref{引理}{11.4}，$D_{S,\allowbreak \rho}$ 为单射，再由指标
\SourcePage{167}
公式可知它可逆。再考虑 $S=D$、$\mu(\rho)=0$ 的情形。\eref{11.13}说明任意非零 $X\in\ker(D_{S,\allowbreak \rho})$ 都是处处非零的函数。这意味着任意两个非零核元素成比例，因为它们的商是 $S$ 上边界取实值的全纯函数，必为常数。因此，在一个边界点求值诱导典范同构 $\ker(D_{S,\allowbreak \rho})\cong\rho(0)$。另一方面，指标公式说明算子为满射，故也有 $\det(D_{S,\allowbreak \rho})\cong\rho(0)$。

\noindent(iv) \emph{一个高维情形。}考虑 $n\geq3$ 时的 $\mathcal L_0\operatorname{Gr}(V)$，并假设 $S=D$。\eref{11.15}的第二个直和因子来自常值环路所成的子空间 $\operatorname{Gr}(V)\subset\mathcal L_0\operatorname{Gr}(V)$。若 $\rho$ 为常值，$(E,\allowbreak F_\rho)$ 显然分解为一维直和分量，(iii) 的论证于是给出 $\det(D_{S,\allowbreak \rho})\cong\lambda^{\mathrm{top}}(\rho(0))$，这说明\eref{11.16}两边限制在 $\operatorname{Gr}(V)$ 上时一致。另一方面，可以构造与 $w_2$ 配对非零的映射 $S^2\to\operatorname{Gr}(V)$，将它视为环路 $S^1\to\Omega\operatorname{Gr}(V)$ 时，所得 $S^1$ 上的行列式线丛非平凡。\bcite{60}第 6 章或\bcite{34}详细进行了这一构造，这里不再重复。这两个观察合起来表明，在 $\mathcal L_0\operatorname{Gr}(V)$ 上有 $w_1(\underline{\det}_S)=T(w_2)+U(\mu)$。

有了这些准备，\eref{11.16}的证明如下。(i) 表明只需考虑 $S=D$。由 (iii)，已知等式在 $\mathcal L_0\operatorname{Gr}(\CC)$ 和 $\mathcal L_{-1}\operatorname{Gr}(\CC)$ 上成立。由 (ii)，若等式在 $\mathcal L_k\operatorname{Gr}(V\oplus\CC)$ 上成立，则它在 $\mathcal L_k\operatorname{Gr}(V)$ 和 $\mathcal L_{k+1}\operatorname{Gr}(V)$ 上也成立；此外，只要 $\dim(V)\geq6$，还可以反向推论。从 (iv) 出发并运用这些观察，便容易得出一般情形。
\end{proof}
\begin{note}\label{st:11.8}
将 Hilbert 空间族 $W^{1,\allowbreak 2}(D,\allowbreak E,\allowbreak F_\rho)$ 平凡化之后——由 Kuiper 定理，这总是可以做到——算子 $D_{D,\allowbreak \rho}$ 在固定 $\rho(0)$ 时定义映射 $\Omega(U_n/O_n)\to\mathfrak F(H,\allowbreak H)$。如\bcite{34}所述，应将它视为 Bott 周期性同伦等价
\[
\mathfrak F(H,H)\simeq\Z\times BO_\infty\simeq\Omega(U_\infty/O_\infty)
\]
的一个近似，其中允许指标相差一个常数平移。
\end{note}

\clearpage
\subsection{Arnol’d 分层}\label{sec:11f}
乘积 $\operatorname{Gr}(V)^2=\operatorname{Gr}(V)\times\operatorname{Gr}(V)$ 按相对位置有一个 Schubert 型分层。每层 $\Sigma^d=\{(\Lambda_0,\allowbreak \Lambda_1):\dim(\Lambda_0\cap\Lambda_1)=d\}$ 上，都有由这些交空间构成的秩 $d$ 向量丛 $\sigma^d$，并带典范同构
\begin{equation}\label{eq:11.17}
(\operatorname{sym}^2\sigma^d)^\vee\cong\nu_{\Sigma^d}.
\end{equation}
可以对 $\operatorname{Gr}(V)^2$ 运用\secref{11b}的路径空间构造，取 $\Sigma=\overline{\Sigma^1}$、$\sigma=\sigma^1$。由此得到局部常值整数函数 $I_Z$，称为\emph{路径的 Maslov 指标}，以及实线丛 $\delta_Z$；它们定义在路径空间 $\mathcal P^2\operatorname{Gr}(V)$ 上，其元素是满足 $s=0,\allowbreak 1$ 时 $\lambda_0(s)\pitchfork\lambda_1(s)$ 的路径 $(\lambda_0,\allowbreak \lambda_1):[0;1]\to\operatorname{Gr}(V)^2$。如同先前
\SourcePage{168}
二次型的情形——事实上该情形可作为局部模型——这里也有涉及 $\Sigma^2$ 的一个容易克服的小困难。

值得把上述构造说得更明确一些。取路径 $(\lambda_0,\allowbreak \lambda_1)$ 和 $s\in[0;1]$。对 $k=0,\allowbreak 1$ 及充分小的 $|r-s|$，选择线性映射族 $\phi_{k,\allowbreak r,\allowbreak s}:\lambda_k(s)\to\lambda_k(r)$，使 $\phi_{k,\allowbreak s,\allowbreak s}=\Id$。\emph{穿越型}是 $\lambda_0(s)\cap\lambda_1(s)$ 上如下定义的二次型：
\[
q_{\lambda_0,\lambda_1}(s):v\longmapsto-\left(\frac d{dr}\right)_{r=s}\omega_V(\phi_{0,r,s}(v),\phi_{1,r,s}(v)).
\]
它与 $\phi_{k,\allowbreak r,\allowbreak s}$ 的选择无关。用\eref{11.17}来表述，若 $(\lambda_0(s),\allowbreak \lambda_1(s))\in\Sigma^d$、$d>0$，则穿越型就是对应于 $(\lambda_0'(s),\allowbreak \lambda_1'(s))$ 法向分量的二次型。对一般位置的路径，空间 $\lambda_0(s)\cap\lambda_1(s)$ 至多一维，而一维交空间上的穿越型非零。\znote{160}{穿越型是 Robbin–Salamon 定义路径 Maslov 指标的出发点。对两族 Lagrangian 子空间 $\lambda_0(s),\lambda_1(s)$，在交点 $s$ 处取两条曲线的速度向量，沿交空间得到一个二次型，其符号记录了交空间在 $s$ 附近如何随参数变化；把路径按穿越点切成小段后，每一段的贡献由该处穿越型的符号给出，Maslov 指标就等于这些符号之和，这与正文用一般机制\eref{11.5}从二次型分层得到的 $I_Z$ 一致。该构造见 Robbin–Salamon《The Maslov index for paths》（Topology 32，1993），nLab 的 Maslov 指标条目列出了该文。[来源：nLab "Maslov index"，https:/\allowbreak /\allowbreak ncatlab.org/\allowbreak nlab/\allowbreak show/\allowbreak Maslov+\allowbreak index（2026-11-15 访问）]}再应用一般机制\eref{11.5}，便得到 $I_Z(\lambda_0,\allowbreak \lambda_1)$ 和 $(\delta_Z)_{\lambda_0,\allowbreak \lambda_1}$。

与\secref{11b}平行，下面说明这些不变量的指标理论含义。选择该处所用的函数 $\psi$。对每个 $(\lambda_0,\allowbreak \lambda_1)\in\mathcal P^2\operatorname{Gr}(V)$，赋予算子 $D_{Z,\allowbreak \lambda_0,\allowbreak \lambda_1}$：它是无穷条带上平凡向量丛 $E=Z\times V$ 的 $\bar\partial$ 算子，边界条件由 Lagrangian 子丛 $F_{\lambda_0,\allowbreak \lambda_1}\subset E|\partial Z$ 给出，其中 $(F_{\lambda_0,\allowbreak \lambda_1})_{(s,\allowbreak k)}=\lambda_k(\psi(s))$。由于端点处的横截条件，它是非退化 Cauchy–Riemann 算子，因而为 Fredholm 算子。于是有
\begin{lemma}\label{st:11.9}
$I_Z(\lambda_0,\allowbreak \lambda_1)=\operatorname{index}(D_{Z,\allowbreak \lambda_0,\allowbreak \lambda_1})$；且 $(\delta_Z)_{\lambda_0,\allowbreak \lambda_1}$ 与 $\det(D_{Z,\allowbreak \lambda_0,\allowbreak \lambda_1})$ 典范同构。
\end{lemma}
\begin{proof}
证明策略与\stref{引理}{11.2}证明的第二种表述相同，故不重复整个论证。这里直接使用粘合公式\eref{11.9}，代替导出\eref{11.8}的初等常微分方程计算。剩下的任务是找出足够多能够显式写出指标与行列式线的例子。一种方法如下；还可采用其他方法，其中有些在技术上稍简单。假设 $V=\R^{2n}=\CC^n$，带标准辛结构和标准复结构 $I_V=i$；对所有 $s$，$\lambda_1(s)=\R^n\times\{0\}$；$\lambda_0(s)$ 是图像 $\{x_2=\alpha(s)x_1\}$，其中 $\alpha(s)$ 为一族对称 $n\times n$ 矩阵。等价地，$\lambda_0(s)=\exp(iA(s))\R^n$，其中
\[
A(s)=\begin{pmatrix}\alpha(s)&0\\0&0\end{pmatrix}.
\]
若对 $(s,\allowbreak t)\in Z$ 应用辛线性映射族 $\Phi(s,\allowbreak t)=\exp((1-t)iA(\psi(s)))$ 改变 $E$ 的平凡化，算子便变为
\begin{equation}\label{eq:11.18}
\begin{aligned}
D_{Z,\lambda_0,\lambda_1}X={}&(\partial_s-J(s,t)iA(\psi(s)))X+J(s,t)\partial_tX\\
&+\Phi(s,t)^{-1}(\partial_s\Phi(s,t))X,
\end{aligned}
\end{equation}
其中 $J(s,\allowbreak t)=\Phi(s,\allowbreak t)^{-1}i\Phi(s,\allowbreak t)$ 是 $\R^{2n}$ 上一族相容复结构。可以将 $J(s,\allowbreak t)$ 连续变形回常值族 $i$；这
\SourcePage{169}
给出\eref{11.18}的一条 Fredholm 算子同伦，故指标与行列式线（在典范同构意义下）不变。同样，\eref{11.18}的最后一项紧支且为零阶项，因而是紧算子，也可通过变形去掉。最终得到
\[
\widetilde D_ZX=(\partial_s+A(\psi(s)))X+i\partial_tX,
\]
以及常值边界条件 $X(s,\allowbreak 0),\allowbreak X(s,\allowbreak 1)\in\R^n\times\{0\}$。椭圆正则性论证说明，每个核元素都可延拓为 $\R^2$ 上的光滑函数，满足 $X(s,\allowbreak t+2)=X(s,\allowbreak t)$、$X(s,\allowbreak -t)=\overline{X}(s,\allowbreak t)$；余核也有类似性质。将定义域与值域限制到这些子空间后，可沿 $t$ 方向作 Fourier 变换，把 $\widetilde D_Z$ 分解为如下单变量微分算子的直和：
\begin{equation}\label{eq:11.19}
\begin{gathered}
d/ds+\alpha(\psi(s)),\qquad
d/ds+\begin{pmatrix}\alpha(\psi(s))&\pi l\cdot\Id\\\pi l\cdot\Id&0\end{pmatrix},\\
l=1,2,\ldots.
\end{gathered}
\end{equation}
\secref{11b}已经研究过这类算子的指标理论。对\eref{11.19}而言，唯一非平凡的贡献来自第一个直和分量，更确切地说，来自使 $\alpha(s)$ 的零空间非平凡的那些 $s$。简单计算表明，这些贡献恰是 $(\lambda_0,\allowbreak \lambda_1)$ 的穿越型。
\end{proof}
\begin{note}\label{st:11.10}
Lagrangian Grassmann 流形的分层由\bcite{2}引入。\stref{引理}{11.9}的第一条是 Floer 的指标定理\bcite{45}。总体而言，本节叙述紧随\bcite{112}、\bcite{113}。上面最后的论证可以视为从 Floer 理论约化到 Morse 理论\bcite{47}的线性版本。
\end{note}
